On fait créer des ondes rectilignes dans la cuve à ondes qui se propagent avec
une vitesse $v=\qty{1}{\v}$, puis on éclaire la surface de l'eau avec un
stroboscope de tel sorte que sa fréquence soit égale à celle des ondes
($f=\qty{10}{\hertz}$), et on voit que tous les points de la surface de l'eau
apparaissent immobiles. On Place deux plaques parallèles dans la cuve de manière
à former une fente de largeur $a$ modifiable. On varie $a$ et on obtient les
deux figures suivantes~:

\begin{figure}[H]
    \centering
    \begin{tikzpicture}[>=Stealth]
        \definecolor{bg}{HTML}{caf0f8}
        \newdimen\h \h=3.5cm
        \newdimen\a \a=1.4cm
        \newdimen\la \la=5mm
        \node[fill,inner sep=0pt,minimum width=5pt,minimum height=\h] (P) {};
        \node[fill=white,inner sep=0pt,minimum width=6pt,minimum height=\a]
            (F) at (P.center) {};
        \draw[<->,shorten <=1pt,shorten >=1pt] (F.south) -- (F.north)
            node[right,pos=0.5] {$a$};
        \foreach \x [count=\i,evaluate=\x] in {1\la,2\la,...\la,6\la}{
            \node[fill=blue,inner sep=0pt,minimum width=1.2pt,minimum height=\h]
                (p\i) at ([xshift=-\x]P.center) {};
            %\node[left color=white,right color=white,middle color=black,
            %      inner sep=0pt,minimum width=5mm,minimum height=\h]
            %    at (\x,0) {};
        }
        \draw[<->,shorten <=0.4pt,shorten >=0.4pt]
            (p3) -- (p4) node[above,pos=0.5] {$\lambda$};
        \foreach \x/\hh [count=\i] in
            {0.7/1.2,1.2/1.6,1.7/2,2.2/2.4,2.7/2.8,3.2/3.2}
        \node[fill=blue,,inner sep=0pt,minimum width=1.2pt,minimum height=\hh cm]
            (p\i) at (\x,0) {};
        \draw[<->,shorten <=0.4pt,shorten >=0.4pt]
            (p3) -- (p4) node[above,pos=0.5] {$\lambda$};
        \draw[rounded corners=8pt,thick] (-3.4,-0.6\h) rectangle (3.6,0.6\h);
        \node[below=3.4mm,font=\footnotesize] at (P.south)
            {Figure~1~: $a=\qty{20}{\cm}$};
        \begin{scope}[xshift=8cm]
            \newdimen\a \a=3mm
            \node[fill,inner sep=0pt,minimum width=5pt,minimum height=\h] (P) {};
            \node[fill=white,inner sep=0pt,minimum width=6pt,minimum height=\a]
                (F) at (P.center) {};
            \foreach \x [count=\i] in {-3,-2.5,...,-0.5}
            \node[fill=blue,inner sep=0pt,minimum width=1.2pt,minimum height=\h]
                (p\i) at (\x,0) {};
            \draw[<->,shorten <=0.4pt,shorten >=0.4pt]
                (p3) -- (p4) node[above,pos=0.5] {$\lambda$};
            \coordinate (S) at ([xshift=3pt]F.center);
            \foreach \r [count=\i] in {0.4,0.9,...,3}
            \draw[line width=1.2pt,blue]
                ([shift=(-40:\r cm)]S) arc (-40:40:\r cm);
            \draw[<->,shorten <=0.6pt,shorten >=0.6pt]
                ([xshift=2cm]F.center) -- ++(0.5,0)
                node[above,pos=0.5] {$\lambda$};
            \draw[rounded corners=8pt,thick] (-3.5,-0.6\h) rectangle (3.5,0.6\h);
            \node[below=3.4mm,font=\footnotesize] at (P.south)
                {Figure~2~: $a=\qty{3}{\cm}$};
        \end{scope}
    \end{tikzpicture}
\end{figure}

\begin{enumerate}
    \item Calculer la longueur d'onde incidente.
    \item Comparer à la largeur $a$ de la fente dans chaque figure.
    \item Décrire, pour chaque figure, ce qui arrive aux ondes lorsqu'elles
          traversent la fente.
    \item L'onde circulaire est appelée l'onde diffractée et le phénomène
          s'appelle phénomène de diffraction. Quelle sont les conditions pour
          que les ondes soient diffractées~?
    \item Comparer la longueur d'onde diffractée avec la longueur de l'onde
          incidente.
\end{enumerate}
